% Figure: a U-tube manometer column oscillating about equilibrium, the
% unsteady-Bernoulli problem; level displacement z(t) decays from an initial
% offset as an undamped (inviscid) sinusoid.
\begin{figure}[htbp]
\centering
\begin{tikzpicture}
\begin{axis}[
  width=12cm, height=6cm,
  xlabel={time $t$ (s)},
  ylabel={column displacement $z$ (m)},
  xmin=0, xmax=8, ymin=-0.6, ymax=0.6,
  axis lines=left,
  grid=major, grid style={enggray!20},
  tick label style={font=\scriptsize},
  label style={font=\small},
  legend style={font=\scriptsize},
  legend pos=north east,
]
% inviscid (undamped) oscillation: z = z0 cos(omega t), omega = sqrt(2g/L)
% L = 1.6 m -> omega = sqrt(2*9.81/1.6) = 3.50 rad/s, period ~1.80 s
\addplot[engblue, very thick, samples=200, domain=0:8]
  {0.5*cos(deg(3.50*x))};
\addlegendentry{inviscid (undamped)}
% lightly damped real column, for contrast
\addplot[engred, dashed, very thick, samples=200, domain=0:8]
  {0.5*exp(-0.18*x)*cos(deg(3.50*x))};
\addlegendentry{real (lightly damped)}
\draw[engblack, thin] (axis cs:0,0) -- (axis cs:8,0);
\node[engblue, font=\scriptsize, anchor=west] at (axis cs:0.2,0.55) {$\omega=\sqrt{2g/L}$};
\end{axis}
\end{tikzpicture}
\caption[Oscillating liquid column in a U-tube]{Displace the liquid in
a U-tube from its level equilibrium and release it: the column
oscillates. The unsteady Bernoulli equation, applied along the
streamline that runs from one free surface through the bend to the
other, turns the time-derivative term $\partial\phi/\partial t$ into
an acceleration of the whole column and yields simple harmonic motion
at $\omega=\sqrt{2g/L}$, with $L$ the total length of the moving
column. The inviscid result is an undamped sinusoid: with no viscosity
the column rings forever. A real column decays, slowly in a wide tube
and quickly in a narrow one, and the decay is the viscous correction
the inviscid analysis leaves out. The frequency, however, is set by
the inviscid balance of gravity against inertia and is accurate well
before the amplitude has decayed appreciably.}
\label{fig:vol03ch11:oscillating-column}
\end{figure}
